\section{Carlos Augusto Preciado Delgado --- Bomberos aeronáuticos}

\subsection{Información básica}
Profesional en sistemas de gestión para la seguridad y salud en el trabajo, técnico profesional en sistemas, docente para bomberos aeronáuticos, entrenador de alturas y labores de alto riesgo. Tiene formación como bombero aeronáutico y coordinador de misiones SAR. Actualmente es coordinador regional y lidera 11 cuerpos de bomberos aeronáuticos. Ha participado en incendios forestales en Yopal, Riohacha y Leticia. No ha participado en incendios de alta montaña.

\subsection{Anotaciones}
\begin{itemize}
  \item Ser cordial y agradecer su disponibilidad.
  \item Aclarar que el interés está en su experiencia como coordinador regional y formador.
  \item Considerar el contexto aeroportuario y de misiones SAR al preparar la entrevista.
  \item No asumir experiencia en incendios de alta montaña.
\end{itemize}

\subsection{Perfil profesional}
\begin{enumerate}
  \pregunta ¿Cuál ha sido su recorrido profesional como bombero aeronáutico, docente y coordinador regional?
  \pregunta ¿Qué responsabilidades tiene actualmente al liderar los 11 cuerpos de bomberos aeronáuticos?
  \pregunta ¿Qué diferencias encontró entre los incendios forestales atendidos en Yopal, Riohacha y Leticia?
  \pregunta ¿Qué aprendizajes de capacitación considera más relevantes para preparar al personal ante incendios forestales?
\end{enumerate}

\subsection{Perspectiva operativa y tecnológica}
\begin{enumerate}
  \pregunta ¿Utilizan actualmente algún software, modelo, mapa o procedimiento para estimar el comportamiento de un incendio? Si es así, ¿qué les aporta y qué limitaciones tiene?
  \pregunta Si una plataforma pudiera simular muchos escenarios posibles de incendios, ¿qué resultados necesitaría mostrar para que fuera útil en planeación o coordinación?
\end{enumerate}


% ---------------------------------------------------------
% CONTENIDO AÑADIDO
% ---------------------------------------------------------

\subsection{Experiencia operativa frente a incendios}

\begin{enumerate}
  \pregunta Cuando se recibe inicialmente el reporte de un incendio forestal, ¿qué información necesita conocer antes de decidir qué recursos movilizar y cómo determina la posible magnitud del incendio durante las primeras etapas de atención?


  \pregunta ¿Cómo se decide cuántas unidades, vehículos, personal u otros recursos deben enviarse inicialmente y cuándo es necesario solicitar apoyo de otros cuerpos de bomberos o entidades?


  \pregunta ¿Qué consecuencias operativas puede tener subestimar o sobreestimar inicialmente la evolución de un incendio y movilizar recursos insuficientes o innecesarios?


  \pregunta Durante los incendios en los que ha participado, ¿qué tipo de información hubiera sido útil tener antes de llegar al lugar?
\end{enumerate}


\subsection{Variables que influyen en la propagación}

\begin{enumerate}
  \pregunta ¿Qué tan relevantes son variables como dirección y velocidad del viento, humedad, vegetación, pendiente y características del terreno?

  \pregunta ¿Cuál de esas variables considera más difícil de conocer o estimar correctamente durante una operación?

  \pregunta ¿Con qué frecuencia cambian esas condiciones durante la atención de un incendio?

  \pregunta Cuando las condiciones cambian, ¿cómo se actualiza la estrategia operativa?

  \pregunta ¿Qué información permitiría anticipar que una estrategia debería modificarse antes de que el cambio sea evidente en campo?
\end{enumerate}


\subsection{Tiempos y necesidades de información}

\begin{enumerate}
  \pregunta ¿Qué decisiones durante un incendio deben tomarse en cuestión de minutos y qué información perdería utilidad si tardara demasiado en llegar?

  \pregunta Si fuera posible conocer una estimación de dónde podría encontrarse el frente del incendio en 30 minutos, una hora o varias horas, ¿para qué decisiones resultaría útil y qué horizonte de tiempo sería más útil desde la perspectiva operativa?

  \pregunta ¿Qué nivel de detalle geográfico necesitaría una estimación para poder utilizarla en campo?

  \pregunta ¿Qué tan importante sería conocer posibles rutas de propagación antes de desplegar personal sobre el terreno?
\end{enumerate}


\subsection{Herramientas actuales y limitaciones}

\begin{enumerate}
  \pregunta ¿Quién suele interpretar actualmente los mapas, datos meteorológicos o información geográfica utilizada durante un incendio, y el personal operativo tiene acceso directo a esa información o debe solicitarla a otras dependencias?


  \pregunta ¿Qué limitaciones tecnológicas encuentran normalmente en campo y con qué frecuencia la falta de conectividad impide utilizar las herramientas?


  \pregunta ¿Qué dispositivos utilizan normalmente para consultar información durante una operación?

  \pregunta ¿Existen herramientas que sean técnicamente buenas pero difíciles de utilizar durante una emergencia?

  \pregunta ¿Considera que las herramientas actuales requieren conocimientos especializados que no todo el personal posee?
\end{enumerate}


\subsection{Simulación y capacitación}

\begin{enumerate}
  \pregunta Desde su experiencia como docente, ¿cómo se enseña actualmente al personal a interpretar y anticipar el comportamiento de un incendio, y qué papel cumplen los ejercicios, simulaciones o escenarios hipotéticos?


  \pregunta ¿Qué características debería tener una simulación para resultar útil como herramienta de entrenamiento?

  \pregunta ¿Sería útil disponer de una biblioteca de diferentes escenarios de incendios para capacitación? ¿Cómo la utilizaría?

  \pregunta ¿Qué errores de interpretación podrían cometer personas con poca experiencia al observar una simulación?

  \pregunta ¿Cómo debería presentarse la incertidumbre de los resultados para evitar que el personal los interprete como una predicción absoluta?
\end{enumerate}


\subsection{Validación de una posible herramienta}

Antes de realizar las siguientes preguntas se explicará brevemente que el proyecto
estudia la posibilidad de disponer de múltiples escenarios de incendios simulados
previamente y, ante una situación real, identificar rápidamente aquellos cuyas
condiciones sean similares a las observadas.

\begin{enumerate}
  \pregunta ¿En qué parte de una operación considera que una herramienta así podría resultar más útil y se utilizaría principalmente durante una emergencia real, para planeación previa, para capacitación o para varias de estas actividades?

  \pregunta ¿Qué variables utilizaría usted para determinar si dos escenarios de incendio son suficientemente similares?

  \pregunta ¿Qué información tendría que mostrar para que un comandante pudiera tomar una decisión apoyándose en ella?

  \pregunta Si la herramienta presentara varios escenarios posibles en lugar de uno solo, ¿esto ayudaría o dificultaría la toma de decisiones?

  \pregunta ¿Cómo deberían compararse esos escenarios para que pudieran interpretarse rápidamente?

  \pregunta ¿Sería importante conocer por qué el sistema seleccionó determinado escenario como similar?

  \pregunta ¿Qué tiempo máximo consideraría aceptable para obtener un resultado?

  \pregunta ¿Sería importante que la herramienta funcionara en condiciones de conectividad limitada?

  \pregunta ¿Qué tipo de capacitación necesitaría el personal antes de utilizar una herramienta de este tipo?

  \pregunta ¿Qué nivel de precisión esperaría para considerar que los resultados son útiles?

  \pregunta En una herramienta de apoyo de este tipo, ¿qué considera más peligroso: que se muestre como amenazada una zona que finalmente no resulte afectada o que no se identifique una zona que posteriormente sí resulte afectada? ¿Por qué?

  \pregunta ¿Utilizaría la herramienta principalmente durante una emergencia real, para planeación previa, para capacitación o para varias de estas actividades?
\end{enumerate}


\subsection{Cierre}

\begin{enumerate}
  \pregunta Si pudiera mejorar una sola capacidad tecnológica disponible actualmente para los cuerpos de bomberos, ¿cuál sería?

  \pregunta ¿Qué característica tendría que demostrar una herramienta como la planteada para que usted considerara confiable utilizarla?

  \pregunta ¿Qué limitación o riesgo considera que nosotros deberíamos tener especialmente en cuenta al diseñarla?

  \pregunta ¿Hay algún aspecto operativo que no hayamos preguntado y que considere fundamental para entender correctamente el problema?

  \pregunta ¿Qué otros bomberos, comandantes, coordinadores, investigadores o entidades considera que sería importante entrevistar?
\end{enumerate}



% =========================================================
% ENTREVISTA 3
% =========================================================

\clearpage
